\documentclass[11pt,a4paper,UTF8]{ctexart}

\usepackage[margin=2.5cm]{geometry}
\usepackage{amsmath,amssymb,bm}
\usepackage{graphicx}
\usepackage{booktabs}
\usepackage{enumitem}
\usepackage{hyperref}
\usepackage{xcolor}
\usepackage{fancyhdr}

\hypersetup{
  colorlinks=true,
  linkcolor=blue,
  urlcolor=blue,
  citecolor=blue
/**
 * @brief: 
 * @note: 
 * @return {*}
 */
}

\pagestyle{fancy}
\fancyhf{}
\lhead{中英文单栏模板}
\rhead{\thepage}
\setlength{\headheight}{14pt}
\renewcommand{\headrulewidth}{0.4pt}

\title{一个适用于中英文的单栏 LaTeX 模板}
\author{Your Name}
\date{\today}

\begin{document}

\maketitle

\begin{abstract}
这是一个可以直接拿来改的单栏模板，适合中文和英文混合写作。你可以把它用于课程论文、简短报告、笔记整理，或者作为更正式文档的起点。

This is a simple single-column template for mixed Chinese and English writing. It is suitable for short reports, class notes, and as a starting point for more formal documents.
\end{abstract}

\tableofcontents
\newpage

\section{引言 Introduction}
这里放引言内容。你可以先交代背景，再说明本文要解决的问题，最后给出结构安排。

For example, you may describe the background first, then state the problem, and finally outline the structure of the document.

\section{正文 Main Text}
这一部分可以写你的主要内容。下面给出几个常见的写法示例。

\subsection{公式示例}
如果需要写公式，可以直接使用行内公式 $E=mc^2$，或者使用独立公式：
\begin{equation}
  \int_0^1 x^2 \, dx = \frac{1}{3}.
\end{equation}

\subsection{列表示例}
\begin{itemize}[leftmargin=2em]
  \item 第一条内容 / First item
  \item 第二条内容 / Second item
  \item 第三条内容 / Third item
\end{itemize}

\subsection{表格示例}
\begin{table}[ht]
  \centering
  \caption{示例表格 Example Table}
  \begin{tabular}{llc}
    \toprule
    项目 Item & 说明 Description & 数值 Value \\
    \midrule
    A & 示例内容 Example text & 1 \\
    B & 示例内容 Example text & 2 \\
    C & 示例内容 Example text & 3 \\
    \bottomrule
  \end{tabular}
\end{table}

\subsection{图表示例}
如果你有图片，可以这样插入：
\begin{figure}[ht]
  \centering
  \fbox{\parbox[c][4cm][c]{0.75\linewidth}{\centering 在这里放图 / Put your figure here}}
  \caption{图示例 Figure Example}
\end{figure}

\section{结论 Conclusion}
这里写结论，总结你在文中得到的主要结果，并可以简单说明后续工作。

In the conclusion, summarize the main results and briefly mention future work if needed.

\section*{致谢 Acknowledgment}
感谢所有帮助过你的人、课程、项目或数据来源。

\end{document}
